\documentclass[11pt]{article}
\usepackage[margin=25mm]{geometry}
\usepackage{fontspec}
\setmainfont{Latin Modern Roman}
\usepackage{amsmath,amssymb,mathtools}
\usepackage{unicode-math}
\setmathfont{Latin Modern Math}
\usepackage{xurl,hyperref}
\hypersetup{colorlinks=true,urlcolor=blue}
\setcounter{secnumdepth}{0}
\setlength{\parindent}{0pt}
\setlength{\parskip}{5pt}
\setlength{\emergencystretch}{3em}
\newcommand{\sref}[2]{\hyperlink{source:#1:#2}{[S#2]}}
\newcommand{\cataloguescope}[1]{\paragraph{Recorded target.}#1}
\begin{document}
\section{2-accessibility of primes}\label{statement}

\subsection{Definitions and mathematical statement}

\paragraph{Shared notation.}$\mathbb N=\{1,2,\ldots\}$, and $\mathbb Z,\mathbb Q,\mathbb R,\mathbb C$ denote the integers, rationals, reals and complex numbers. Any other conventions are specified locally.


Let $\mathcal P=\{2,3,5,\ldots\}$ be the positive primes. For $D\subset\mathbb N$, a $D$-diffsequence of length $k$ is a sequence of positive integers $x_1<\cdots<x_k$ with $x_{i+1}-x_i\in D$ for each $i$. The set $D$ is two-accessible if every coloring $c:\mathbb N\to\{0,1\}$ has monochromatic $D$-diffsequences of every finite length. The successive gaps need not be equal.
\paragraph{Statement recorded in the supplied source.}
Is $\mathcal P$ two-accessible? Equivalently, for every $c:\mathbb N\to\{0,1\}$ and every integer $k\geq1$, must there exist $x_1<\cdots<x_k$ with all $c(x_i)$ equal and every difference $x_{i+1}-x_i$ prime? \sref{statement}{2}
\subsection{Short English statement}
Does every two-coloring of the positive integers have arbitrarily long monochromatic increasing sequences whose successive gaps are prime?
\subsection{Sources}
\begin{description}
\item[\hypertarget{source:statement:1}{S1}] ULAM.AI and the UnsolvedMath Contributors, \emph{UnsolvedMath: A Curated Collection of Open Mathematics Problems}, catalogue entry. CC BY 4.0. \url{https://huggingface.co/datasets/ulamai/UnsolvedMath}.
\item[\hypertarget{source:statement:2}{S2}] Oscar Quester, \emph{The Primes are 2-Accessible}, arXiv:2606.00410v1 (2026), Definitions 1--2 and Theorem 2. The theorem gives an affirmative answer and proves that the accessibility degree of the primes is exactly two. \url{https://arxiv.org/abs/2606.00410}.
\end{description}
\label{end:statement}
\end{document}
